\documentclass[11pt]{article}

\usepackage[margin=1in]{geometry}
\usepackage[utf8]{inputenc}
\usepackage[T1]{fontenc}
\usepackage{lmodern}
\usepackage{microtype}
\usepackage{amsmath, amssymb}
\usepackage{graphicx}
\usepackage{booktabs}
\usepackage{xcolor}
\usepackage{hyperref}
\usepackage[round]{natbib}

\hypersetup{
  colorlinks=true,
  linkcolor=blue!50!black,
  citecolor=blue!50!black,
  urlcolor=blue!50!black,
}

\title{Persona Prompting and Sycophancy Compound:\\
       Evidence from a Production Evidence-Grounded Synthesis System}

\author{Jacob\thanks{Code, data, and per-run telemetry available at
        \url{https://github.com/onedusk/cce}. Reproduction notes in
        \texttt{paper/README.md} of the repository.}}

\date{April 30, 2026}

\begin{document}
\maketitle

\begin{abstract}
% TODO: 150--250 words.
% Frame: persona prompting and sycophancy share a root cause (RLHF preference
% optimization overriding factual grounding) and should compound when combined.
% State the experimental result and the novel contribution (implied uncited
% claims via rhetorical structure).
\end{abstract}

\section{Introduction}
\label{sec:intro}

Three findings about large language models have accumulated independently
in the recent literature. First, expert persona prompting that improves
alignment, tone, and format compliance also degrades factual accuracy on
discriminative tasks: \citet{hu_prism_2026} report all persona variants
reduced MMLU accuracy from a 71.6\% baseline to a 68.0\% best, and
\citet{basil_pretend_2025} replicate the effect on GPQA Diamond and
MMLU-Pro across six models. Second, production language models exhibit
systematic sycophancy: \citet{cheng_sycophantic_2025} test 11 frontier
models in two preregistered experiments (N=1{,}604) and find AI affirms
users 50\% more than human comparators, and 80\% or more on inputs the
authors classified as clearly wrong, against 40\% for humans.
\citet{shapira_rlhf_2026} provide a mathematical account of the
amplification mechanism: RLHF optimization against a learned reward
mechanically biases policy outputs toward agreement when human preference
data carries any agreement signal. Third, instruction-following operates
as compositional coordination of specialized skills with finite
representational capacity, with verification concentrated at generation
boundaries \citep{rocchetti_instructions_2026}.

Each of these effects is documented in isolation. Their compounding
interaction is not. The compositional-coordination account predicts that
adding a persona prompt should consume monitoring bandwidth that would
otherwise support factual grounding, and that the consumption should be
worse under additional resource pressure. The RLHF-amplification account
predicts that when monitoring capacity is consumed, the surface-compliance
signal that the reward model has learned to value should dominate over
factual accuracy. The two predictions converge: a sycophancy-prone model
given an expert persona, especially under word-budget compression that
creates conflicting constraints, should exhibit accuracy degradation
beyond either effect alone.

We test the converged prediction in a production
evidence-grounded synthesis system, the Content Curation Engine (CCE).
CCE was built for citation-backed content production, not as an
experimental apparatus, but the variables required for this study are
already present in its production telemetry: a per-path persona prompt
field that can be toggled on or off, per-unit metrics for confidence,
citation coverage, source diversity, claim count, and training-data
leakage rate, and a fixed evidence store that holds source material
constant across persona conditions. Section~\ref{sec:cce} describes the
apparatus.

% TODO: contribution-claims paragraph — fill after the persona on/off
% experiment (task #6) lands. Will state: (1) magnitude and direction of
% the persona-on vs. persona-off delta on leakage; (2) the compounding
% interaction with word-budget tier; (3) the implied-uncited-claim rate
% reported in §7 as an undocumented category of persona-induced accuracy
% loss; (4) the spectrum-principle mitigation that closes the gap.

\section{Related Work}
\label{sec:related}

\paragraph{Persona prompting and accuracy.}
\citet{hu_prism_2026} (PRISM) document that every expert persona variant
they tested reduced MMLU accuracy below the 71.6\% baseline, with the best
variant landing at 68.0\%. They attribute the loss to the model becoming
reluctant to express uncertainty under persona instructions and instead
confabulating confidently. PRISM proposes a gated LoRA adapter that
applies the persona only when downstream task type is a generative one.
\citet{basil_pretend_2025} replicate the persona-degradation effect on
GPQA Diamond and MMLU-Pro across six models, finding persona prompts
generally did not improve accuracy and sometimes reduced it. The two
results together rule out the possibility that the effect is an artifact
of MMLU specifically.

\paragraph{Sycophancy in production systems.}
\citet{cheng_sycophantic_2025} measure sycophancy across 11 frontier
LLMs using two preregistered experiments with 1{,}604 participants.
Models affirm user actions 50\% more often than human comparators; on
inputs classified as clearly wrong, models agree 80\% or more of the time
against 40\% for humans. Users rated sycophantic responses as higher
quality and reported greater trust, creating a reinforcing feedback loop
in any system that uses user feedback as a training signal.
\citet{chandra_spiraling_2026} formalize the downstream user-cognition
risk: across 10{,}000 simulated conversations, even Bayesian-optimal
users developed high confidence in false beliefs from sycophantic
chatbots, and the effect persisted when false claims were eliminated and
users were warned about sycophancy. The OpenAI GPT-4o incident of April-May
2025 is the production-scale instance: a model update combining user
feedback signals, persistent memory, and fresher training data compounded
into behavior that validated delusions and urged impulsive actions, with
at least 14 reported deaths and 5 lawsuits before the model was retired
\citep{openai_gpt4o_sycophancy_2025}.

\paragraph{Mechanism.}
\citet{shapira_rlhf_2026} provide a mathematical account of the
amplification mechanism. RLHF optimization against a learned reward
mechanically biases policy outputs toward agreement when human preference
data carries any agreement signal; the bias direction is characterized by
the covariance between agreement features and reward features under the
base policy. \citet{wang_overridden_2025} complement this with a
mechanistic-interpretability analysis showing sycophancy is a structural
override of learned knowledge in deeper layers, not a surface artifact;
they also report that first-person opinion statements (``I believe...'')
induce higher sycophancy than third-person framings, while user-expertise
framing has negligible impact. \citet{rocchetti_instructions_2026} provide
converging evidence that instruction-following is not a single
constraint-satisfaction mechanism but compositional coordination of
specialized skills with finite capacity. Through diagnostic probing across
nine tasks in three models, they find general probes underperform
task-specific specialists, cross-task transfer is weak, causal ablation
reveals sparse asymmetric dependencies rather than shared representations,
and constraint-satisfaction operates as dynamic monitoring during
generation with a verification spike at the EOS token. Their account
predicts that adding a persona constraint should consume coordination
bandwidth that would otherwise support factual grounding skills.

\paragraph{Measurement under constraint.}
\citet{baxi_compliance_2025} build a benchmark that independently measures
constraint compliance and semantic accuracy under output-length
compression. The two dimensions are statistically near-orthogonal
(Pearson r = 0.193). Removing RLHF helpfulness signals improved constraint
compliance by 598\%, identifying RLHF-trained helpfulness as the dominant
cause of constraint violation under resource pressure.
\citet{sun_length_2025} test 30 LLMs under varying output-length budgets
and report a U-curve: too-short and too-long budgets both degrade
performance, with the optimal model size and prompt style varying by
budget. \citet{hong_sycon_2025} introduce SYCON Bench for multi-turn
sycophancy, finding that alignment tuning amplifies sycophantic
behavior while model scaling and reasoning optimization reduce it; a
third-person persona designed to resist user pressure increased
sycophancy resistance by up to 63.8\% in debate scenarios.
\citet{qiu_rhetorical_2025} provide a methodology for separating
rhetorical style from content via counterfactual generation, which we
build on in Section~\ref{sec:implied} when measuring contrastive-frame
prevalence independent of the topical content of the contrast.

\section{Hypothesis}
\label{sec:hypothesis}

We hypothesize that persona prompting and sycophancy share a root cause
(RLHF preference optimization that values surface compliance over factual
grounding) and that they should compound rather than add when combined.
The qualitative form of the prediction follows from
\citet{rocchetti_instructions_2026}: if instruction-following is
coordination of specialized skills with finite capacity and verification
concentrated at generation boundaries, a persona constraint introduces a
persistent coordination burden that competes with factual-grounding
skills throughout generation. Under additional resource pressure (a tight
word budget that forces the writer to prioritize among required
sections), the coordination should fail in a direction consistent with
the bias trained into the reward model. \citet{shapira_rlhf_2026}
characterize that direction: when the reward model has learned an
agreement-is-good heuristic from human preference data, the policy's
optimization gradient favors stance-affirming outputs over truth-tracking
ones at the margin.

Concretely, with a sycophancy-prone base model, an expert persona prompt
that suppresses uncertainty expression, and a word-budget tier that
creates contention between required-section coverage and citation density,
we expect (i) higher training-data leakage rate when persona is on than
when persona is stripped; (ii) the leakage gap to widen at tighter word
budgets, with persona-on confidence holding steady or improving while
leakage rises (the surface-compliance signal eats accuracy); and
(iii) higher density of contrastive frames carrying implied uncited
claims (Section~\ref{sec:implied}) when persona is on, because the
persona instructions reward the rhetorical posture the frames produce.
In the language of \citet{baxi_compliance_2025}, persona prompting should
move output toward the constraint-compliant, semantic-accuracy-degraded
quadrant of the joint distribution.

% TODO: quantitative pre-registered predictions — derive specific
% effect-size point estimates from the Shapira covariance characterization
% applied to CCE's measured agreement-feature distribution. Worked through
% in task #2 and committed (with timestamp) before the experiment runs.

\section{The Content Curation Engine: An Evidence-Grounded Synthesis System}
\label{sec:cce}

\subsection{Pipeline architecture}

The Content Curation Engine is a production system for citation-backed
content synthesis. The pipeline flow is: a curation request invokes a
Discoverer that retrieves candidate sources subject to a configurable
source policy (domain allow/deny rules, reputation weighting, recency
constraints); each retrieved source is parsed into evidence objects with
full provenance (source URL, hash, excerpt offsets) and stored in an
EvidenceStore; a Writer produces a draft conditioned on the stored
evidence and the path-specific configuration; a Verifier checks every
claim in the draft against the evidence store; a quality gate routes the
draft to PASS (publish), FAIL (rewrite, up to a per-risk-profile
iteration cap of 2-4), or REVIEW (human moderation).

The pipeline enforces a single foundational invariant: \emph{no citation,
no ship}. The Writer's prompt restricts generation to claims grounded in
stored evidence objects, which the Writer cites by ID. The Verifier reads
each cited claim against the corresponding evidence excerpt and assigns
a status (supported, unsupported, uncited, or leakage, where leakage is
defined as a claim with no corresponding evidence object that appears
to derive from training data). The gate computes coverage (fraction of
claims with positive Verifier status) and confidence (the citation-density
weighted aggregate) per content unit, and refuses to publish below
threshold. The invariant is what makes CCE useful as an experimental
apparatus for studying accuracy: leakage is measured continuously in
production, not as a synthetic eval.

The closest related system in the published literature is STORM
\citep{shao_storm_2024}, a long-form Wikipedia-style writing system that
also retrieves grounded sources and synthesizes article-length output.
STORM concentrates on the pre-writing stage, modelling research as
multi-perspective question generation against simulated topic experts and
producing an outline before any prose is written. The authors note
``source bias transfer and over-association of unrelated facts'' as
problems they identified but did not solve. CCE differs in scope and
emphasis: where STORM optimizes the breadth and organization of
pre-writing, CCE operates the full writer-verifier-gate loop in
production with the trust-contract invariant enforced at the gate, and
treats the over-association problem (in our terminology, implied claims
carried by rhetorical structure) as a verifiable category rather than an
acknowledged limitation. Section~\ref{sec:implied} extends STORM's
observation by quantifying one specific form of the over-association
pattern in production output.

\subsection{Per-unit telemetry}

Each content unit produced by CCE carries the following measurements,
recorded per generation iteration with full lineage:
\textit{confidence} (citation-density weighted verification score),
\textit{coverage} (claim pass rate),
\textit{source\_diversity} (independent-source count weighted by reputation),
\textit{citation\_count} and \textit{claim\_count} (raw counts),
\textit{leakage} (count and rate of training-data hallucinations
detected), and per-stage timing. Iteration-level deltas allow direct
observation of how the writer-verifier loop moves the unit through the
quality space across rewrites.

A separate humanization layer, opt-in via the engine configuration,
captures style-level metrics: sentence-length variance, lexical
diversity, suppressed-vocabulary count, formulaic-transition density,
contrastive-frame density, hedging frequency, and em-dash density per
1000 words. These ride on the same per-iteration record and are logged
as \texttt{JobStage.SCORE} entries for threshold calibration. The scorer
is purely programmatic; an Editor agent applies stylistic rewrites
conditional on the scorer's pass/fail; a downstream implied-claim
checker (described in Section~\ref{sec:implied}) extends the trust
contract from explicit to implied claims.

\subsection{The persona variable}

Each content path (Learn, Explore, Apply in the current taxonomy) carries
a \texttt{prompt\_addendum} field that applies a path-distinct persona
instruction to the Writer's system prompt. The current production values
specify ``calm authority,'' ``intellectual credibility,'' and ``thoughtful,
warm,'' respectively. The field is an ordinary part of the path
configuration: it can be set, cleared, or varied without modifying engine
code, which makes it directly available as the independent variable for
the persona on/off experiment.

\subsection{Operational baseline}

For grounding the apparatus, an early autonomous run completed with 66
cited claims drawn from 15 independent sources, a 98.6\% verification
score, and zero training-data fabrications detected. The first draft of
that run produced two weak claims that the gate flagged; the
writer-verifier loop produced a rewrite that cleared the gate without
human intervention. We use this to set expectations for the
persona-stripped condition in later sections: at the operating point of
the current production system, the engine is capable of zero-leakage
output, so observed leakage in the persona-on condition is attributable
to the persona prompt and word-budget tier rather than to a generally
broken pipeline.

\section{Experimental Design}
\label{sec:design}
% Persona on/off across topics and path types (Learn/Explore/Apply), word-budget
% tier crossing. Dependent variables: confidence, coverage, leakage,
% contrastive-frame density, implied-uncited-claim count. Statistical tests.
% Pre-registered predictions from Section~\ref{sec:hypothesis}.

\section{Results}
\label{sec:results}
% Persona-on vs.\ persona-off deltas, broken out by path and word-budget tier.
% Test pre-registered predictions. Headline tables / figures:
%   - persona effect on leakage
%   - compounding interaction with resource pressure
%   - implied-uncited-claim count delta

\section{Implied Claims via Rhetorical Structure}
\label{sec:implied}

A citation-grounded synthesis pipeline that checks every explicit claim
against its evidence can still produce confidently false output when some
claims are encoded implicitly in rhetorical structure. This section
identifies one such category, measures its prevalence in the CCE corpus, and
argues that closing the gap requires extending the trust contract from
explicit to implied claims.

\subsection{The pattern}
\label{sec:implied-pattern}

A contrastive frame asserts a position by dismissing an alternative.
Constructions of the form ``\textit{X is not A. It is B}'' or ``\textit{Unlike
X, Y does Z}'' are common in persuasive prose and frequent in
LLM output \citep{qiu_rhetorical_2025}. Each instance carries two claims: a
positive claim about the writer's preferred concept, and a negative claim
about the dismissed alternative. The positive claim is typically backed by an
explicit citation in our corpus. The negative claim usually is not, because
the writer treats the dismissed alternative as obvious rather than as a claim
requiring evidence.

Consider a sleep-treatment example from the CCE evidence store: ``Unlike
sleeping pills, CBT-I addresses the underlying causes of sleep problems
rather than just relieving symptoms.'' The positive claim about CBT-I can be
verified against guideline literature in the store. The negative claim
about sleeping pills---that they only relieve symptoms and do not address
underlying causes---is asserted without citation, and is contradicted by other
evidence in the same store: orexin receptor antagonists, for instance, target
wakefulness pathways rather than sedating around them. The verifier did
not flag the contrast because the negative claim was never stated as a
verifiable proposition; it was carried by the syntax of the contrast.

This generalizes. Any contrastive frame has a positive pole (the writer's
argument) and a negative pole (what is being dismissed). Verifiers built
around explicit claim extraction see only the first.

\subsection{Empirical prevalence in the CCE corpus}
\label{sec:implied-empirical}

To estimate how often this pattern appears in production output, we ran two
regular expressions matching the canonical syntactic shapes (``\texttt{is not
X. It is Y}'' and structural variants including ``\texttt{are not X. They are
Y}'' and the ``\texttt{not just X. but also Y}'' scope-expansion) across CCE
runs covering 11 distinct topics across the \texttt{learn} and
\texttt{explore} content paths. The regex returned 46 matches.

Each match was hand-labeled with a three-way rubric:
\textit{parasitic} (the dismissed pole is rhetorical scaffolding; collapsing
the sentence to ``\textit{It is B}'' loses no factual content),
\textit{factual} (the dismissed pole carries an independent claim; collapse
would destroy information; the canonical class here is the safety-relevant
counterexample ``the drug is not dangerous; it is essential''), and
\textit{ambiguous} (either reading is defensible without wider context).

The audit returned 44 parasitic matches (95.7\%), 0 factual matches, and 2
ambiguous matches (4.3\%). Treating the zero observed factual cases as a
binomial sample, the 95\% one-sided Clopper-Pearson upper bound on the
factual-match rate is approximately 6.3\%. Treating the ambiguous cases as
worst-case factual, the upper bound rises to approximately 14\%. Both bounds
indicate that the dominant function of contrastive frames in this
production corpus is rhetorical reframing of a single referent, not
genuine alternative-structure between two factually distinct claims.

The two ambiguous cases are worth describing because they bound the
class. Match \#10 (``This is not the tiredness that sleep repairs. It is a
deeper depletion'') carries a clinical distinction between burnout exhaustion
and ordinary tiredness that mechanical collapse would strip. Match \#37
(``The response is not broken. It is simply being applied to situations it
was never built for'') carries an independent reassurance beyond what the
positive clause conveys. Neither case crosses into the safety-relevant
factual class the rubric was guarding against.

\subsection{The spectrum principle and verifier extension}
\label{sec:implied-mitigation}

A surface-level mitigation that suppresses contrastive frames at the
vocabulary or syntactic level addresses the AI-detectability concern but
leaves the underlying binary reasoning pattern intact: the writer finds new
words to express the same dichotomy. A more durable mitigation reshapes the
reasoning step that produces the binary in the first place. We refer to this
as the \emph{spectrum principle}: when introducing a distinction between two
concepts, the writer must include at least one condition under which the
disfavored side is valid or relevant. The evidence store typically contains
material supporting both sides of any real contrast, so this constraint is
also a better use of the evidence already collected.

Applied to the sleep example, the spectrum-aware version reads: ``Sleeping
pills can break the cycle of acute insomnia quickly, but they do not change
the habits and thought patterns that keep the problem coming back. CBT-I
targets those directly, which is why guidelines now recommend it as the
first-line treatment, though for severe cases short-term medication
alongside CBT-I can outperform either alone.'' Same evidence, no implied
uncited claim, and incidentally less detectable as AI-generated because the
context-dependent treatment of the contrast resembles human domain
expertise.

The verifier extension is mechanical. When the syntactic patterns above
are detected, the verifier extracts the implied negative claim about the
dismissed side and searches the evidence store for supporting material. If
supporting evidence is found, the frame is flagged as containing an uncited
implied claim, and the writer receives a rewrite directive that names the
specific evidence objects supporting the dismissed side and the context in
which the dismissed side is valid. The cost is one additional retrieval per
matched frame. In the audit corpus, this would have flagged the 44
parasitic matches for spectrum rewrites and either flagged the 2 ambiguous
matches for human review or, in the worst-case binding, inserted up to 14\%
spurious flags.

\subsection{Connection to the trust contract}
\label{sec:implied-contract}

The pipeline's foundational invariant is that no claim ships without a
citation. Implied claims violate that invariant invisibly: the citation
verifier sees a sentence whose explicit propositions are all backed,
and certifies it. The accuracy gap is structural, not stylistic, and
addressing it is properly scoped to the citation verifier rather than the
humanization layer. The style benefit (less detectable as AI-generated) follows from the
accuracy fix as a side effect. This framing aligns with the
constraint-compliance versus
semantic-accuracy decomposition of \citet{baxi_compliance_2025}: implied
claims are a class of constraint violation that present as compliant under
surface checks while degrading semantic accuracy.

\section{Mitigation}
\label{sec:mitigation}

The mitigation literature on AI-generated text divides cleanly into
surface fixes (vocabulary suppression, burstiness scoring, em-dash
density limits, formulaic-transition blocklists) and structural fixes
that change the reasoning step that produces a class of pattern. The
distinction matters here because surface fixes leave the underlying
accuracy problem intact: a writer that cannot use the word ``delves''
will find a synonym for the same rhetorical move, and a writer that
cannot use ``Unlike X, Y'' will find a different syntactic shape for
the same binary construction. The implied-claim audit in
Section~\ref{sec:implied} is direct evidence: the parasitic frames are
not failures of vocabulary control; they are failures of binary-versus-
spectrum reasoning.

\subsection{Surface mitigations: useful but bounded}

The vocabulary literature is now precise enough to operationalize.
\citet{qiu_rhetorical_2025} provide a counterfactual methodology for
isolating rhetorical style from content; the corpus studies aggregated
in our internal review identified 379 excess style words elevated since
the ChatGPT release, with ``delves'' at 28$\times$ expected frequency,
``underscores'' at 13.8$\times$, and ``showcasing'' at 10.7$\times$, and
a focal-word set of 21 high-signal entries (delve, intricate, nuanced,
underscore, showcasing, comprehensive, multifaceted, pivotal,
commendable, meticulous, realm, landscape, tapestry, crucial, robust,
plus six others). CCE's humanization scorer maintains these as
configurable blocklists in YAML so the markers can track the
coevolution of detector and model: a known signal degrades in detection
value as soon as it becomes known, because models suppress it under
pressure. Burstiness (sentence-length variance), lexical diversity
(type-token ratio), and structural-template detection are similarly
useful programmatic checks, but each one detects an effect downstream
of the reasoning step that produces the AI fingerprint, not the
reasoning step itself.

\subsection{Resource-pressure findings from the CCE corpus}

A length-tier study run on a single topic (stress physiology) with the
same evidence store and the same source policy across three word-budget
configurations isolates the effect of resource pressure on the leakage
rate. Tier B (Learn 2000 / Explore 1500 / Apply 1000) is the only
configuration where all three paths cleared the gate; the Learn path
went from 4 leakage claims on iteration 1 to 0 on iteration 2 with a
1.000 verification score. Tier A (Learn 1500 / Explore 1200 /
Apply 800) entered a whack-a-mole pattern on the Explore path: claim
counts rose from 46 to 48 to 51 across three iterations as the writer
added claims to address verifier feedback, with 2-3 leakage events on
each iteration despite the rewrite. Tier C (Learn 3000 / Explore 2500 /
Apply 1500) failed by a different mechanism: claim volume gave the
verifier more surface area, dropping per-iteration confidence below
threshold even when the unsupported rate per claim was low. The
takeaway is that leakage in this pipeline is a function of word budget
relative to required-section count, not just of model behavior, and
that the failure mode at the tight end is precisely the
training-data-substitution behavior the implied-claim and persona
mechanisms above predict.

\subsection{Structural mitigation: spectrum principle and implied-claim
verification}

The structural counterpart to vocabulary suppression is the spectrum
principle described in Section~\ref{sec:implied}: the writer prompt
adds a constraint that distinctions between two concepts must include
at least one condition under which the disfavored side is valid, and
the verifier extends the trust contract by extracting the implied
negative claim from contrastive frames and searching the evidence store
for material supporting the dismissed pole. Where surface mitigations
produce a writer that finds new words for the same binary, the spectrum
principle changes the underlying reasoning: the writer has to represent
the relationship between the two ideas before producing prose. The
2026-04-17 calibration sweep over 36 archival drafts set the soft
thresholds for the humanization gate at the values currently shipped in
\texttt{config/humanization\_markers.yaml}; the em-dash threshold
(4.0/1000) is an editorial target rather than an engine floor, since
em-dash overuse is the AI fingerprint we want the editor to address
rather than refuse on.

\subsection{Coevolution and configuration-as-policy}

The coevolutionary dynamic from \citet{qiu_rhetorical_2025} and earlier
work has a direct implication for engineering: any specific marker list
has a finite shelf life, and the scorer's configuration is itself a
moving target. CCE addresses this by keeping the marker lists, density
thresholds, and rhetorical-pattern regexes in operator-editable YAML
that ships separately from the engine code, so the scorer can be
retuned without redeployment. This is the same pattern the citation
literature recommends for hallucination evals: the test set decays as
the model learns its boundaries, so the test-set-as-config decision
matters more than the specific test cases on it.

\section{Discussion}
\label{sec:discussion}

\subsection{Generalization beyond content curation}

The pipeline structure on which the experimental result rests is
domain-pluggable. Strip the content-specific framing and the apparatus
is a query-driven evidence-grounded synthesis loop: discover sources,
extract facts with provenance, synthesize a grounded output, verify
every claim against the evidence, gate on a trust threshold. The
``no citation, no ship'' invariant is the load-bearing piece, and it
is domain-agnostic. Adapter implementations exist or are sketched for
web pages (current), API-backed sources (SEC EDGAR, financial data
providers), code repositories (docs, issues, version-pinned API
references), document corpora (PDF parsing for filings, patents, case
law), and structured data sources. Trust heuristics that hardcode
academic conventions (.gov and .edu domains, peer-review status) become
pluggable reputation strategies per domain: jurisdiction authority and
court level for legal research, regulatory-filing precedence over
analyst notes for financial research, official-doc and version-recency
weights for programming research. Output formats parameterize from
prose to comparison tables, structured briefs, code snippets with
sourced explanation, or jurisdictional summaries. The persona-leakage
finding generalizes wherever a stylistic prompt addendum is applied to
a writer that operates against a fixed evidence store.

\subsection{Contradiction as a feature}

In the current content domain, contradiction between sources is treated
as a problem the gate must resolve: the writer picks a side, the
verifier checks the cited side, and the disagreement is flattened. In
research-flavored domains, the disagreement is often the value. A
competitive-intelligence query that returns ``Source A reports market
share of X, Source B reports Y, here is the methodology and recency for
each'' is more useful than one that picks a winner. The same
infrastructure that powers implied-claim detection in
Section~\ref{sec:implied} (extract negative pole, search evidence store
for support) is the natural primitive for surfacing contradictions as
first-class outputs: the spectrum principle at the architectural level.
\citet{shao_storm_2024}'s multi-perspective question-asking at the
pre-writing stage moves in a similar direction; lifting that primitive
from outline construction into the verifier-gate path would close the
loop. This direction also shifts the gate from binary PASS/FAIL to
confidence-weighted continuous output, where each claim carries a
confidence score and a source-count, and the consumer applies a
threshold appropriate to their use (a regulatory submission requires
0.95+; a casual reader accepts 0.7+).

\subsection{Agentic grounding}

A speculative direction worth naming: the same pipeline can operate as
a trust layer for autonomous agents that take external action. Instead
of ``write me an article about X,'' the call shape becomes ``before you
recommend Y, verify it against current sources.'' The agent's action
routes through the evidence loop before it executes anything that
depends on external facts. This is plausibly the highest-value
application of the pipeline beyond content, and it is the place where
the persona-and-sycophancy compounding finding has the largest
operational stakes: an agent that gets persona instructions to be
``helpful'' or ``confident,'' and that operates against a sycophancy-
prone base model, will compound the failure mode in a domain where the
output is an action rather than a paragraph.

\subsection{Limitations}

The experimental result reported here is from a single production
synthesis system on a content domain, with a fixed model family and a
fixed evidence store. The compositional-coordination account of
\citet{rocchetti_instructions_2026} is itself a finding from three
small models (Llama 3.1 8B, Gemma 2 2B, Qwen 2.5 0.5B); whether the
finite-capacity prediction holds at frontier scale is an open question.
The persona-on/off experimental design (Section~\ref{sec:design}) is
sufficient for testing the directional prediction but does not
disentangle the contributions of the persona instruction, the path
structure, and the word-budget tier individually beyond their
two-way interactions. Replication on a second synthesis pipeline (a
research-domain instance built on the same architecture but a different
evidence-source mix) would strengthen the generalization claim.

% TODO: expand the limitations subsection after the persona on/off
% experiment lands. Items to address concretely once we have data:
% (a) external validity beyond the wellness-content corpus;
% (b) effect-size sensitivity to model version;
% (c) whether the implied-claim detection pipeline produces false
%     positives at a rate that would change the gate's operating point.

\section{Conclusion}
\label{sec:conclusion}

% \nocite{*} forces all bib entries to render, useful while the body is stubbed.
% Remove this line once real \citep{...} calls populate the sections.
\nocite{*}

\bibliographystyle{plainnat}
\bibliography{references}

\end{document}
